\label{sec:lift}
\subsection{The essence of Zhou's paper (after Kuan)}
Zhou studies a random surface growth in the quarter plane: interlacing particles that push and block, with a reflecting wall at $0$. It
is a one-parameter family of Plancherel measures for $O(\infty)$, so its characters are of type B/D. Three results matter here.
\begin{enumerate}[leftmargin=*]
\item \emph{The process is determinantal} along space-like paths, through a generalized $L$-ensemble.
\item \emph{It is the shadow of a free non-commutative walk.} On $U(\mathfrak{so}_{N+1})$ take the states $\langle X\rangle_\kappa=D(X)\kappa(I)$, with
$\kappa_t(O)=e^{t\Tr(O-I)}$. Their moments are an exponential of a trace,
\[\langle F_{i_1j_1}\cdots F_{i_mj_m}\rangle_t=\sum_\pi t^{|\pi|}\prod_B\Tr\prod_{b\in B}F_{i_bj_b},\]
which is Faà di Bruno again. Then
$P_t=(\mathrm{id}\otimes\langle\cdot\rangle_{\kappa_t})\circ\Delta$ is a semigroup that preserves the centre. Through the Harish-Chandra isomorphism, central
elements act by symmetric polynomials in the shifted labels, i.e.\ the particle positions. Restricted to the centre, $P_{t/2}$ is the
particle process of one level (Theorem~4.2). Restricted to the Gelfand--Tsetlin subalgebra, generated by the centres of the nested chain
$\mathfrak{so}_2\subset\dots\subset\mathfrak{so}_{N+1}$, it is the whole multi-level process (Theorem~4.3).
\item \emph{Gaussian free field, and the wall's type matters.} Moments of central elements converge to a three-dimensional Gaussian free
field. That field differs from the one of overlapping Wishart spectra, although both live to the right of a wall at $0$.
\end{enumerate}
The lessons are these.
\begin{itemize}[leftmargin=*]
\item[(L1)] A process with complicated classical interactions (blocking, pushing, reflection) is the commutative shadow of a free
non-commutative process. One computes upstairs.
\item[(L2)] The shadow is taken on a commutative subalgebra built from a nested chain.
\item[(L3)] What kind of wall it is (reflecting, or a hard edge) decides the fluctuations.
\end{itemize}
The mechanism is the one of Section~\ref{sec:hopf}: $P_t=(\mathrm{id}\otimes\langle\cdot\rangle_t)\Delta$ and $\phi_\lambda=(\lambda\otimes\mathrm{id})\Delta$ coincide for
cocommutative $\Delta$. A dynamics is convolution with a semigroup of functionals.

\subsection{The network's lift: the quasi-free field and its gate shadow}
For the network the upstairs object is the quasi-free (Gaussian) field of pre-activations, the backbone $(\mu_l,C_l)$.
\begin{itemize}[leftmargin=*]
\item The commutative subalgebra is generated by the gates, i.e.\ the cell indicators. This is the diagonal of stage~11's tile
groupoid.
\item These algebras increase with depth, matching $\mathrm{Gap}_l\subset\mathrm{Gap}_{l+1}$ in stage~14. That is a nested chain as in (L2).
\item The gate gas is the restriction of the field's state to it.
\end{itemize}
\begin{proposition}[the fold; proved]\label{prop:fold}
$\relu=\frac12(\mathrm{id}+|\cdot|)$.
\begin{enumerate}[leftmargin=*]
\item The identity half is linear. It commutes with the $O(n)$ symmetry of the weight ensemble, sends quasi-free laws to quasi-free laws,
and creates no cumulants (Proposition~\ref{prop:layeraction}).
\item The fold $|\cdot|$ is the quotient by the sign group $(\Z/2)^n$. Together with the relabelling symmetry $S_n$ of hidden units, this is
the hyperoctahedral group $B_n$, the Weyl group of type~B.
\begin{itemize}
\item In the layer's own frame $z=Wx$ the walls are the reflection hyperplanes of $(\Z/2)^n$, and the cells are its chambers.
\item So the network alternates $O(n)$-equivariant linear maps with $B_n$-equivariant reflecting walls.
\item All births come from the folds.
\end{itemize}
\item At a fold, the even and odd Hermite sectors are Laguerre sectors with half-integer parameters:
\begin{itemize}
\item $\He_{2k}(x)=(-2)^kk!\,L_k^{(-1/2)}(x^2/2)$ and $\He_{2k+1}(x)=(-2)^kk!\,x\,L_k^{(1/2)}(x^2/2)$.
\item The parameters $a=\pm\frac12$, which index Zhou's two level parities through the Jacobi weights $(1-x)^a(1+x)^{-1/2}$, are the
spectral signature of a reflecting wall. Our fold carries the same pair in its Laguerre limit.
\end{itemize}
\item The folded law is an image sum, $p_{|z|}(y)=p(y-\mu)+p(y+\mu)$ for $y>0$. This is the method of images of the reflecting walk.
\end{enumerate}
\end{proposition}

\subsection{The gate gas: arcsine pairs and stars}
\begin{proposition}[proved to leading order; checked]\label{prop:stars}
Let $z\sim N(0,R)$ have unit variances and correlations $\rho_{ij}$, and let $\varepsilon_i=\operatorname{sign}z_i$. Then
$\E\varepsilon_i\varepsilon_j=\frac2\pi\arcsin\rho_{ij}$, and the fourth joint cumulant is
\[
\kappa(\varepsilon_1,\varepsilon_2,\varepsilon_3,\varepsilon_4)=-\frac4{\pi^2}\sum_{i=1}^4\prod_{j\ne i}\rho_{ij}+O(\rho^4).
\]
\end{proposition}
\begin{proof}
Write $\operatorname{sign}=\sum_{k\ \rm odd}a_k\He_k/k!$ with $a_k=2\He_{k-1}(0)\varphi(0)$, so $a_1=\sqrt{2/\pi}$ and $a_3=-\sqrt{2/\pi}$. In the
Gaussian diagram expansion the two-vertex components reproduce the arcsine pairings. The lowest connected diagram with odd degrees on
four vertices is the star of degrees $(3,1,1,1)$, of weight $a_3a_1^3=-(2/\pi)^2$.
\end{proof}
\emph{Check} (orthant probabilities by Genz quadrature): the ratio of the exact cumulant to the star formula is $1.013$, $1.038$, $1.074$
at correlation scales $0.1$, $0.2$, $0.3$. Two readings follow.
\begin{itemize}[leftmargin=*]
\item \emph{Rounding.} The arcsine law is Grothendieck--Krivine rounding of the linear, quasi-free correlation. Tsirelson's theorem
realizes the linear correlations themselves with Clifford observables. The classical gates pay the rounding, up to the Grothendieck
constant in correlation strength.
\item \emph{Stars are the non-quasi-free content.} What the gate gas has beyond pairings begins with stars, gates with three legs. These
are the hub-rooted stars of stage~9 (P2) and the births of the chain. Stage~12 called the memory the ``non-quasi-free remainder''; this
is its first diagram.
\end{itemize}

\subsection{Two sectors, and where the missing information lives}
Every statistic an estimator forms from the instance splits into two parts.
\begin{itemize}[leftmargin=*]
\item An $O(n)$-invariant part: spectral and Gram data, traces, kernels. This is the isotropic sector, and it contains the modular one.
\item A $B_n$-specific remainder, which depends on the neuron basis: Hadamard and diagonal readouts, hub pairings.
\end{itemize}
For the invariant sector, (L1)--(L2) apply in their usual form. Linear spectral statistics of nested rank-one (hub) chains have Gaussian
fluctuations whose covariance is computable from the ensemble, as for Wishart minors. For the $B_n$ sector no such reduction exists.
Corollary~\ref{cor:nofreelunch} says that this is where the information an estimator lacks lives. Lesson (L3) appears as the split
$\relu=\frac12(\mathrm{id}+|\cdot|)$. The identity half carries the first jet, so it transports errors; the fold half creates them.
